\section{Experimental setup}\label{sec:setup}

\paragraph{Datasets}
We use four standard benchmarks for fair node classification (Table~\ref{tab:data}). \emph{German} (credit risk; sensitive attribute: gender) and \emph{Credit} (credit default; age) are similarity graphs over individuals, \emph{Bail} (bail decisions; race) links defendants with similar records, and \emph{Pokec-z} (a social network; region) is the regional subset of Pokec used in~\citep{dai2021say}. The sensitive attribute is removed from the node features. Each run draws a stratified 50/25/25 train/validation/test split and partitions the nodes over $K=5$ clients with a Dirichlet($\alpha=0.3$) split over labels, which produces strongly non-IID clients. Each client keeps the subgraph induced by its nodes, so cross-client edges are lost (between 21\% and 38\% of the edges survive, Table~\ref{tab:data}). The split, the partition and the model initialisation all depend on the seed.

\begin{table}[t]
\centering
\caption{Datasets after preprocessing. Edges are undirected. ``Kept'' is the fraction of edges inside a client after the $K=5$ Dirichlet partition (seed 0). $\ncap$ is the public cap on a client's per-group validation count.}\label{tab:data}
\begin{tabular}{lrrrcccr}
\toprule
Dataset & Nodes & Edges & Features & $P(y{=}1)$ & $P(s{=}1)$ & Kept & $\ncap$ \\
\midrule
German  & 1{,}000  & 21{,}742    & 26  & 0.700 & 0.310 & 0.21 & 50 \\
Bail    & 18{,}876 & 311{,}870   & 16  & 0.376 & 0.506 & 0.38 & 1{,}000 \\
Credit  & 30{,}000 & 1{,}421{,}858 & 12  & 0.779 & 0.090 & 0.34 & 2{,}750 \\
Pokec-z & 67{,}796 & 617{,}958   & 276 & 0.081 & 0.352 & 0.29 & 4{,}050 \\
\bottomrule
\end{tabular}
\end{table}

\paragraph{Protocol constants}
All methods train for $T=100$ rounds with $E=3$ full-batch local epochs, AdamW with learning rate $0.01$ and weight decay $10^{-5}$, two-layer backbones with 64 hidden units (GAT: 4 heads) and dropout $0.3$. For \clamp{} we use $C=1$, $\eta=0.5$, fairness budget $\tau=0.03$, statistics privacy $\varepsilon=8$ at $\delta=10^{-5}$ for the attribute-level guarantee of Theorem~\ref{thm:dp}, and release period $P=1$, except $P=10$ on German, whose validation groups are too small for per-round releases (Section~\ref{sec:privacy}). We set $\ncap$ to the largest per-client validation group size rounded up to a multiple of 50. In our experiments this cap is computed from the partition. In deployment it would be fixed from public dataset sizes before training, and Theorem~\ref{thm:dp} assumes it is public. \clamp{} uses a GAT backbone; \clamp{}-GCN swaps in GCN.

\paragraph{Baselines}
We compare with: FedAvg with GCN and GAT backbones~\citep{mcmahan2017communication}; the fair GNNs FairGNN~\citep{dai2021say} and FairSIN~\citep{yang2024fairsin} trained with FedAvg; the fair FL methods FairFed~\citep{ezzeldin2023fairfed} and FedFACT~\citep{fedfact2025} (closed-form group mean-matching post-processing); the federated graph methods F$^2$GNN~\citep{f2gnn2023}, FairGFL~\citep{fairgfl2026} (its aggregation rule) and FedGraph-Fair~\citep{fedgraphfair2026} (its distributionally robust reweighting, without personalisation); and the private fair methods PUFFLE~\citep{corbucci2024puffle} and FedFDP~\citep{ling2024feddp}, both at gradient privacy $\varepsilon=8$. For FedFDP we clip at the batch level rather than per sample, as in our DP-FedAvg implementation, because per-node gradients are not independent in a message-passing GNN. FedFDP optimises fairness of performance across clients, not group fairness, so it is not expected to reduce the demographic-parity gap. \clamp{}-ClipMean is \clamp{} with $\eta=0$: the same clipped secure aggregation without the fairness controller. It isolates the effect of the multiplier.

\paragraph{Metrics and statistics}
We report test AUC, the soft demographic-parity gap $\dpd$ of Eq.~\eqref{eq:dpd} on test nodes, and the equal-opportunity gap $\eod$ at threshold $0.5$. For \clamp{} we also report the final multiplier and the certificate $\lambda^{(R)}/(\eta R)$ of Proposition~\ref{prop:cert}. Each configuration is run with ten seeds. We test each baseline against \clamp{} with a two-sided Wilcoxon signed-rank test paired by seed, with Holm correction across baselines within each dataset and metric.

\paragraph{Attacks}
Update attacks (Section~\ref{sec:byz}) use $K=10$ clients of which $f\in\{1,2,3,4\}$ are Byzantine: sign flipping, scaling, ALIE~\citep{baruch2019little}, inner-product manipulation~\citep{xie2019fall} and a fairness-poisoning attack whose local loss rewards a larger demographic-parity gap while still fitting the labels. Under \clamp{} attacker updates are clipped to the norm bound $C$ like all others, so they always pass the validity check. We compare \clamp{} with FedAvg, the coordinate median, the trimmed mean and Krum, all with a GAT backbone; the last three need individual updates and cannot run under secure aggregation. Statistics attacks (Section~\ref{sec:statattack}) let the Byzantine clients train honestly but falsify their statistics with $N_a=\ncap$ and either swapped group rates (\emph{swap}) or both rates equal to the honest pooled rate (\emph{mid}), which drives the reported gap toward zero.

\paragraph{Compute}
Runs used an NVIDIA T4 GPU (Credit, Pokec-z and the attack, sweep and ablation studies) and a 4-core CPU (German, Bail and the privacy study). The code, configuration and raw results of every run are released with the paper.
